\section{CIFER interfaces and query-local compilation}

\subsection{Finite observation interfaces}

The same ideal mathematical object may support many finite interfaces.  A continuous real function, for example, can be consumed through pointwise values, interval enclosures, a modulus on a compact interval, a finite graph approximation, or polynomial coefficients.  A universal representation is rarely the simplest route to \iPRA.  \CIFER\ therefore makes the interface relative to a query.

\begin{definition}[Finite observation interface]
A finite observation interface for inputs \(x\) and demand \(q\) consists of:
\begin{enumerate}
\item an evidence code \(e\);
\item a quantifier-free verifier \(V(x,q,e)\);
\item an observation code \(o=\operatorname{obs}(e)\);
\item a correctness formula \(C(x,q,o)\);
\item an embedded derivation
\[
 V(x,q,e)\vdash_{\PRAAx} C(x,q,\operatorname{obs}(e)).
\]
\end{enumerate}
The interface may additionally contain a primitive-recursive producer \(E(x,q)\) and a proof of \(V(x,q,E(x,q))\).  If the producer is unavailable, the existence or validity formula becomes a residual boundary.
\end{definition}

The observation may itself be relational.  In that case \(e\) codes a finite list \(L\), the verifier proves that every member is valid and that the list is nonempty or complete for the finite query, and the consumer works with all members or with a stable observable.

\subsection{Demand transformers}

Every interface declares what input information is needed for an output demand.  For a deterministic operation \(f\), a demand transformer is a primitive-recursive function
\[
 d_f(q,\alpha)=p,
\]
where \(q\) is the requested output precision and \(\alpha\) collects finite side information such as denominator exponents, coefficient heights, or a supplied gap.  The local correctness theorem has the form
\[
  C_X(x,p,a)\to C_Y(f(x),q,F(a)).
\]

For a graph node with several inputs, the transformer returns one demand per predecessor.  The reverse pass takes componentwise maxima over all consumers.  When a transformer's side information depends on demands elsewhere, the graph gives a monotone system rather than a simple acyclic recurrence.  Section~\ref{sec:demand-closure} solves that system under an explicit cap.

\subsection{Evidence classes}

The framework uses a small number of evidence patterns repeatedly.

\paragraph{Exact calculation evidence.}
An exact node computes a primitive-recursive term and proves its defining equation.  Examples are integer polynomial evaluation, rational arithmetic, finite sums, matrix multiplication in fixed dimension, and list filtering.

\paragraph{Trace evidence.}
A trace node codes an initial state, a transition, a fuel, and an invariant.  Its public result is extracted from the terminal state.  Bounded search, Euclidean algorithms, descent, dynamic programming, finite Gaussian elimination, and interval iteration all fit this pattern.

\paragraph{Certificate evidence.}
A certificate need not be generated by the verifier.  Examples include a factorization, a Bezout identity, an ideal-combination witness, a nonzero resultant, a graph coloring, a finite cover, or a rational point.  The verifier is primitive recursive and the theorem says what follows if it accepts.

\paragraph{Relational evidence.}
A relational node returns a finite cloud of candidates together with validity and compatibility tables.  It supports collision-prone or noncanonical data.

\paragraph{Observation evidence.}
An observation node proves that all valid candidates have a common label, lie in a small interval, induce the same homology map, or yield the same bounded consumer result.

\paragraph{Boundary evidence.}
A boundary node states a producer result that is used downstream but not constructed in the current package.  Examples are ``this code is a complete height bound,'' ``these are rational covering representatives,'' and ``this rational is an effective zero-free constant.''

\subsection{Proof dependencies versus data dependencies}

It is useful to distinguish three graphs.

\begin{enumerate}
\item The \emph{proof graph} is a finite DAG.  A local proof may use only earlier nodes and the permanent context.
\item The \emph{data-flow graph} records which computed codes feed which operations.  It is usually a DAG after bounded loops are packaged as trace nodes.
\item The \emph{demand graph} records how requested resources propagate.  It may contain cycles, because a bound can depend on a precision that depends on the same bound.
\end{enumerate}

Conflating these graphs leads to confusion.  A demand cycle does not make the proof circular if it is resolved by a finite fixed-point computation before the proof run.  Conversely, a topologically ordered executable computation does not justify a proof node unless the local object derivation is supplied.

\subsection{Query specialization}

A query consists of:
\begin{itemize}
\item a target formula \(G\);
\item permanent input formulas \(\Gamma\);
\item a requested finite observation or precision;
\item an application library of possible interfaces.
\end{itemize}
The host compiler selects only the interfaces reachable from \(G\), specializes their terms and formulas, computes residual boundaries and demands, and emits a concrete proof schedule.

This specialization serves three purposes.

First, it limits the amount of formal arithmetic.  A theorem about a polynomial of fixed degree needs only fixed-dimensional tuple operations, not an internal general-purpose algebra system.  Second, it makes resource bounds sharper: the demand analysis sees the actual consumer.  Third, it avoids universal evaluation.  The generated object program is a fixed composition of \code{PR} constructors.

\subsection{A schema for deterministic contracts}

Let \(P(x)\) be a precondition, \(f(x)\) a fixed primitive-recursive output term, and \(Q(x,y)\) a postcondition.  A deterministic contract consists of
\[
 P(x)\vdash Q(x,f(x)).
\]
Composition is direct cut.  If
\[
 P\vdash Q(f),\qquad Q(y)\vdash R(g(y)),
\]
then instantiate the second local derivation at \(f\) and cut the first proof into it.  On a finite DAG, repeated simultaneous cut composes all nodes.

A deterministic node is appropriate only when the output really is a total primitive-recursive function of the declared input.  It should not be used to hide an unbounded search or a discontinuous selector.

\subsection{A schema for relational contracts}

A relational contract has a finite output list \(L(x,q)\) and a verifier proving
\[
 \operatorname{Nonempty}(L)\wedge
 \forall i<|L|\;R(x,q,L_i).
\]
If completeness is needed, add
\[
 \forall y\;\bigl(R^*(x,q,y)\to \exists i<|L|\;E(y,L_i,q)\bigr),
\]
where the quantification over ideal objects is replaced in the actual application by its finite coded or observational form.  The consumer may:
\begin{enumerate}
\item search the finite list for an additional decidable property;
\item propagate all candidates to another finite relation;
\item prune candidates by backward viability;
\item apply a stable observer;
\item require a gap and collapse to an exact label.
\end{enumerate}

The relation is therefore compositional.  Nonuniqueness is not a defect to be repaired prematurely.

\subsection{A schema for conditional producers}

Suppose \(B(x,b)\) says that \(b\) is a valid producer output and the downstream finite theorem is
\[
 B(x,b)\vdash G(x,b).
\]
There are three legitimate package levels.

\begin{enumerate}
\item \textbf{Verifier level:} prove the implication for arbitrary supplied \(b\).
\item \textbf{Producer level:} define a primitive-recursive \(p(x)\) and prove \(B(x,p(x))\).
\item \textbf{Closed level:} cut the producer proof into the verifier proof.
\end{enumerate}

When only the first is available, \CIFER\ reports \(B\) as a residual.  It does not introduce a noncomputable host constant \(b\), mark it opaque, and then present the result as a theorem of \iPRA.

\subsection{Residual normalization levels}

There are several notions of ``minimal boundary,'' and they should not be confused.

\begin{description}[style=nextline]
\item[Reachability normalization.] Remove boundary nodes not used by the target.
\item[Duplicate normalization.] Merge repeated occurrences of the same boundary node.
\item[Dependency expansion.] Replace any derived pseudo-boundary by the primitive unresolved leaves needed to derive it.
\item[Declared dominance.] If the library contains an embedded proof that one boundary follows from others, use it as an ordinary local node and eliminate the dominated boundary.
\item[Logical minimization.] Find a smallest logically sufficient subset of arbitrary formulas.  This is not a computable generic service and is not claimed.
\end{description}

The canonical \CIFER\ residual is minimal only in the first three syntactic senses, plus any explicitly supplied dominance proofs.  This level is deterministic, auditable, and sufficient for modular work.

\subsection{Proof-carrying interface records}

At the Lean level, an interface should not be a structure containing only semantic fields.  It should contain, or generate, a local embedded derivation.  Schematically:

\begin{lstlisting}
structure LocalNode (Gamma : List (Fm PTm n)) where
  formula : Fm PTm n
  deps : List (Fm PTm n)
  proof : deps ++ Gamma |- [PRAAx] formula
\end{lstlisting}

Executable metadata may be stored alongside the proof, but the proof field is decisive.  An application library can then expose two APIs:
\begin{enumerate}
\item a high-level builder that computes syntax and evidence;
\item a small trusted-facing record whose final field has type \code{Prf PRAAx}.
\end{enumerate}

This is the proof-theoretic analogue of proof-carrying code: the consumer checks the derivation rather than trusting the generator.

\subsection{Migration rule for existing packets}

An existing proof packet can be migrated incrementally.

\begin{enumerate}
\item Identify its headline target formula.
\item List all inputs currently described by phrases such as ``choose,'' ``there exists an effective,'' ``for a sufficiently large,'' or ``given a complete list.''
\item Classify each as a finite producer, a boundary, a relation, or a gap.
\item Turn every executable loop into a trace node with one explicit invariant.
\item Turn native Lean helper theorems into specifications for local object derivations, not into assumed bridges.
\item Compile the downstream DAG and record the residual context.
\item Close one small benchmark end to end before extending the domain library.
\end{enumerate}

This strategy preserves existing mathematical work while making proof status progressively more precise.
